\section{Correspondence with the thesis: first proposal}\label{sec:thesis}
The thesis full text has now been consulted for the following anchors. Printed
page numbers refer to the thesis, not its archive cover. This targeted reading
locates obligations; it is not a complete audit of the thesis's semantics.

\begin{center}\small
\begin{tabularx}{\linewidth}{@{}p{28mm}p{15mm}X@{}}\toprule
Thesis anchor & Pages & Proposed role\\\midrule
5.4.10 & 101 & Context transition-system semantics; identify a flat fragment.\\
6.2.3--6.2.4 & 116 & HML syntax and satisfaction.\\
6.2.5--6.2.6 & 117 & Tests with truth polarity and oracle satisfaction.\\
6.2.7--6.2.8 & 118 & Exhaustive tests and full-agreement setting; no finite prototype claim.\\
8.3.1 & 160--162 & Expanded execution patterns; restrict data-variable scope.\\
8.3.3--8.3.4 & 162--164 & Pattern validation and positive/negative substitutions.\\
8.4.1--8.4.2 & 165--167 & Behavioral/coordination annotations with object identity.\\
8.4.5 & 169--170 & Annotated pattern validation.\\
8.4.7--8.4.8 & 171--172 & Reduced substitution families and selected test sets.\\\bottomrule
\end{tabularx}
\end{center}

\paragraph{Interpretation issue requiring author review.}
The printed rules in Definition 8.3.3, p.~163, place the diamond continuation at
source $q$ and include an existential failing-continuation rule in its negative
relation. These details were checked visually in the PDF. Literal use is not the
standard nondeterministic HML semantics specified in Section~\ref{sec:contract}
and the preceding section: the continuation is evaluated at $q'$, and failure of
one successor does not refute a diamond. We propose to state standard HML
explicitly, then record an agreed interpretation or erratum before asserting
correspondence with that definition. The annotated rules require the same review.
This is an open interpretation issue, not a claimed correction accepted by the author.

\paragraph{Subuniformity bridge.}
Thesis annotations identify behavioral/coordination formulas and object instances.
An equation-unfolding branch in the prototype is not itself such an annotation.
Definition 8.4.7 constrains \emph{sets} of annotated substitutions, including
consistency of chosen values under matching annotations; Definition 8.4.8 chooses
from reduced families. We must decide whether the algorithm is intended to produce
one member of such a family, cover feasible annotation classes, or satisfy a
stronger equality. Its current one-witness-per-computed-region policy does not
resolve this choice.

\begin{obligation}[P5, initial restricted target: per-test preservation]
For a declared thesis fragment and a translation $\tau$ of its states, labels and
finite formulas, establish initial-state correspondence and truth preservation
for each admitted ground test:
\[
 S_{\mathrm{th}},q_0\models\varphi
 \quad\Longleftrightarrow\quad
 S_{\mathrm{impl}},m_0\models\tau(\varphi).
\]
The interpretation of the thesis relation must first be agreed.
\end{obligation}
One sufficient route is a relation matching accepted steps in both directions and
preserving the relevant observation predicates, followed by induction on finite
HML. A one-way simulation is insufficient for unrestricted negation and boxes.
This target deliberately precedes equality of reduced test sets, representative
adequacy or bounded completeness. Ground evaluator/SMT agreement, integer bounds,
partiality and constructor equality remain explicit premises to discharge.

\todo{Agree the reading of Definitions 8.3.3 and 8.4.5; define the flat source
fragment and translations; decide the reduction-family coverage target; choose the
initial algebra subset. No theorem relating computed regions to thesis
subuniformity classes is established in this draft.}
